\documentclass[11pt,a4paper]{article}

\usepackage[T1]{fontenc}
\usepackage[utf8]{inputenc}
\usepackage{lmodern}
\usepackage[margin=2.6cm]{geometry}
\usepackage{amsmath,amssymb}
\usepackage{graphicx}
\usepackage{booktabs}
\usepackage{caption}
\usepackage[colorlinks=true,linkcolor=black,citecolor=black,urlcolor=NavyBlue]{hyperref}
\usepackage[usenames,dvipsnames]{xcolor}
\usepackage{titlesec}
\usepackage{abstract}
\usepackage{enumitem}
\usepackage{microtype}

\captionsetup{font=small,labelfont=bf,width=.95\textwidth}
\titleformat{\section}{\normalfont\large\bfseries}{\thesection}{0.7em}{}
\titleformat{\subsection}{\normalfont\normalsize\bfseries}{\thesubsection}{0.7em}{}
\setlength{\parskip}{0.35em}
\setlength{\parindent}{0pt}
\renewcommand{\abstractnamefont}{\normalfont\bfseries\small}
\renewcommand{\abstracttextfont}{\normalfont\small}

\title{\vspace{-1.2cm}\bfseries The Walk-Forward Bundle:\\[2pt]
\large Specification-Curve Analysis for Trading-Strategy Validation}

\author{Francesco Landolfi\\[2pt]
\small Epiphany Alpha\\
\small \texttt{francesco.landolfi@epiphany-alpha.com}}

\date{\small August 2026 \quad$\cdot$\quad Working paper, not peer reviewed}

\begin{document}
\maketitle

\begin{abstract}
\noindent A walk-forward back-test is widely treated as the honest form of
historical evaluation: parameters are fitted on a past window and evaluated on
the window that follows, so no future information enters the fit. We argue that
this discipline, while necessary, is not sufficient. The procedure itself
requires a set of construction choices --- in-sample window length, re-selection
cadence, ranking objective, correlation cap --- for which no principled answer
exists, and each choice yields a different out-of-sample equity curve. Reporting
one such curve constitutes an unreported search over construction space.

We propose reporting the \emph{bundle}: run the full grid of defensible
walk-forward constructions, retain every out-of-sample path, and report the
resulting distribution rather than any element of it. This is the
multiverse-analysis and specification-curve methodology of the empirical
social sciences, transposed to strategy validation.

We demonstrate the procedure on two datasets processed identically. On a
synthetic price series containing no edge by construction, the median of 432
walk-forward constructions is an annualized out-of-sample Sharpe of $-0.27$,
while the best is $+0.18$ --- a number that, reported alone, would be
indistinguishable from a discovery. On S\&P~500 daily data evaluated
out-of-sample from 2000, the same 432 constructions give a median of $+0.33$
with every construction positive and a dispersion roughly half as wide. The
diagnostic is not the level but the shape: the gap between the best
construction and the median falls from $0.45$ to $0.10$. We also report the
computational work required to make exhaustive evaluation cheap enough to be
routine, in which a per-window memoization of the ranking metric accounted for
essentially all of a $3.8\times$ speedup. A reference implementation is released
under the MIT licence.
\end{abstract}

\vspace{0.4em}

\section{The construction problem}

Consider a researcher with a family of candidate strategies --- parameter
variants of one model, several models, or one model across several assets ---
and a walk-forward protocol for selecting among them. On each pass, candidates
are ranked on an in-sample window, a subset is retained, and that subset is
traded over the period immediately following. The out-of-sample legs are
concatenated into a single track record.

Nothing in that description is unsound. The resulting equity curve contains no
look-ahead: the selection at each fold uses only data preceding the period it
is judged on. This is a real improvement over a single in-sample optimisation
and is standard practice.

The difficulty lies one level above. Before the protocol can be run, the
researcher must fix:

\begin{itemize}[leftmargin=1.3em,itemsep=1pt,topsep=3pt]
  \item the in-sample window length --- 12 months? 24? 48?
  \item whether that window expands from a fixed origin or rolls;
  \item the re-selection cadence, which also fixes the length of each
        out-of-sample leg;
  \item the objective on which candidates are ranked --- Sharpe, Calmar,
        Sortino, total return;
  \item how many candidates are retained, and how correlated they may be.
\end{itemize}

These are not nuisance parameters with known correct values. They are choices,
and each combination produces a different out-of-sample curve. When a single
curve is reported, a search over construction space has taken place and has not
been disclosed --- typically not even to the researcher, who experienced the
process as a sequence of reasonable decisions rather than as an optimisation.
The out-of-sample discipline internal to each run offers no protection, because
the selection that matters occurred above it.

This is precisely the structure Gelman and Loken termed the garden of forking
paths \cite{gelman2013}, and the response developed in the empirical social
sciences after the replication crisis is instructive: rather than defend one
specification, report all of them. Steegen et al.\ \cite{steegen2016} call this
a multiverse analysis; Simonsohn et al.\ \cite{simonsohn2020} formalise the
reporting as a specification curve. In finance the closely related concern ---
that the reported performance of a strategy reflects the number of trials that
produced it --- is treated by White's reality check \cite{white2000}, by the
multiple-testing haircut of Harvey et al.\ \cite{harvey2016}, and by the
backtest-overfitting programme of Bailey et al.\
\cite{bailey2014,bailey2015}. Our contribution is neither a new statistic nor a
new correction. It is the observation that the walk-forward protocol has its own
multiverse, that this multiverse is small enough to enumerate exhaustively, and
that its shape is directly interpretable.

\section{Method}

\subsection{Setup}

Let $R \in \mathbb{R}^{T \times k}$ be a matrix of candidate return streams,
column $j$ being the realised P\&L of candidate $j$ at time $t$, and let
$V \in \mathbb{R}^{T \times k}$ be the corresponding traded notional
(turnover). We require that $R$ be \emph{causal}: each column must be the
output of a rule whose parameters were not fitted using data after $t$. Nothing
downstream can repair a violation of this requirement.

A construction $c$ is a tuple
\[
c = \big(L,\; s,\; a,\; m,\; n,\; \rho,\; K\big),
\]
with $L$ the in-sample window length in months, $s$ the re-selection cadence
(equal to the out-of-sample leg length), $a \in \{\text{expanding},
\text{rolling}\}$, $m$ the ranking objective, $n$ the shortlist size, $\rho$ the
correlation cap and $K$ the maximum number of candidates deployed.

\subsection{One construction}

Given a first out-of-sample date $t_0$, the schedule generates windows
$[\,u_i, v_i\,]$ with $v_i = t_0 + i\cdot s$ and $u_i = v_0 - L$ when $a$ is
expanding, $u_i = v_i - L$ when rolling. The schedule terminates $s$ months
before the end of the sample, so that every selection is followed by a period
on which it can be judged.

For each window, in order:

\begin{enumerate}[leftmargin=1.5em,itemsep=1pt,topsep=3pt]
  \item \textbf{Rank.} Evaluate $m$ on $R[u_i{:}v_i]$ column-wise and retain the
        best $n$ candidates.
  \item \textbf{De-correlate.} Traverse the shortlist best-first, admitting a
        candidate only if its correlation with every already-admitted candidate
        is below $\rho$; stop at $K$ admissions. This is deliberately greedy and
        deterministic: explicability matters more than optimality when the
        object of the exercise is a robustness claim.
  \item \textbf{Gate on economics (optional).} Discard candidates whose realised
        P\&L per unit of turnover, $\sum_t R_{tj} / \sum_t V_{tj}$, falls below
        a floor. This is the quantity that must exceed transaction cost.
  \item \textbf{Evaluate.} Equal-weight the survivors over
        $(v_i, \, v_i + s]$ and record the resulting return and turnover series.
\end{enumerate}

The out-of-sample leg begins strictly after the in-sample window closes; the two
never share an observation. Concatenating the legs yields a single continuous
out-of-sample return stream $r^{(c)}$, from which we compute an annualized
Sharpe ratio $\widehat{S}(c)$.

\subsection{The bundle}

Let $\mathcal{C}$ be the Cartesian product of the admissible values of each
construction parameter. The object we report is the collection
\[
\mathcal{B} \;=\; \big\{\, r^{(c)} \;:\; c \in \mathcal{C} \,\big\},
\]
together with the induced distribution of $\widehat{S}(c)$ over $\mathcal{C}$.
The summary statistics of interest are the median, the dispersion, the fraction
of constructions with positive out-of-sample Sharpe, and --- the quantity we
find most diagnostic --- the gap
\[
\Delta \;=\; \max_{c \in \mathcal{C}} \widehat{S}(c) \;-\;
              \operatorname{med}_{c \in \mathcal{C}} \widehat{S}(c),
\]
which measures how much apparent performance is available to a researcher who
selects a construction after seeing the results.

A grid should span choices the researcher genuinely cannot rule out, and no
others. Adding axes on which one has no opinion inflates $|\mathcal{C}|$ without
adding information; adding axes on which one setting is known to be wrong
imports a known-bad configuration into the reported distribution.

\section{Two experiments}

Both experiments use an identical pipeline: a sweep of $105$ trend-following
rules (Donchian breakout, exponential-moving-average crossover, and time-series
momentum, across six fast and six slow look-backs), volatility-targeted, charged
2 basis points per unit of traded notional, executed one bar after signal. The
construction grid is
$L \in \{12,24,36,48\}$, $s \in \{3,6,12\}$, $a \in \{\text{expanding},
\text{rolling}\}$, $m \in \{\text{Sharpe}, \text{Calmar}, \text{Sortino}\}$,
$\rho \in \{0.3,0.5,0.7\}$, $K \in \{5,10\}$, with $n = 20$ throughout ---
$|\mathcal{C}| = 432$.

\subsection{A synthetic null}

The first price series is generated: a volatility-clustering, fat-tailed process
with constant drift. It has the texture of a market --- long quiet stretches
punctuated by violent moves, so trend rules behave as they do on real data ---
and, by construction, contains nothing a look-back rule can predict. The
individual candidates have a median annualized Sharpe of $-0.09$: after costs,
the family loses money, as it should.

Figure~\ref{fig:syn} shows all 432 out-of-sample paths. The median construction
returns an annualized out-of-sample Sharpe of $-0.27$ and only $5\%$ of
constructions end positive. The best returns $+0.18$.

\begin{figure}[t]
\centering
\includegraphics[width=\textwidth]{fig/fig1_synthetic.pdf}
\caption{All 432 walk-forward constructions on the synthetic null. Each faint
line is one construction's concatenated out-of-sample return stream. Every line
is a genuine walk-forward result: fitted on a past window, traded on the window
after it, no look-ahead. The data contains no edge.}
\label{fig:syn}
\end{figure}

That $+0.18$ deserves emphasis. It is arithmetically correct, strictly
out-of-sample, produced by a defensible construction, and entirely an artefact
of having tried 432 of them on data engineered to contain nothing. A researcher
who ran this grid, selected the best construction and reported its curve would
not be committing fraud; they would be doing what the field's reporting
conventions permit.

\subsection{S\&P 500}

The second series is the S\&P~500 daily index from 1927, with the first
out-of-sample date set to January 2000. The same 105 rules and the same 432
constructions are applied without modification. Individual candidates have a
median annualized Sharpe of $0.45$.

\begin{table}[t]
\centering
\small
\caption{The bundle on both datasets. Annualized out-of-sample Sharpe over 432
walk-forward constructions.}
\label{tab:main}
\begin{tabular}{lrr}
\toprule
& \textbf{Synthetic} & \textbf{S\&P 500} \\
\midrule
Median                                   & $-0.27$ & $+0.33$ \\
Minimum                                  & $-0.79$ & $+0.06$ \\
Maximum                                  & $+0.18$ & $+0.43$ \\
Dispersion (s.d.)                        & $0.16$  & $0.09$  \\
Constructions with $\widehat{S} > 0$     & $5\%$   & $100\%$ \\
$\Delta$ (best $-$ median)               & $0.45$  & $0.10$  \\
\bottomrule
\end{tabular}
\end{table}

\begin{figure}[t]
\centering
\includegraphics[width=\textwidth]{fig/fig2_comparison.pdf}
\caption{The same procedure on both datasets. (a) The synthetic null: a wide
fan, mostly below zero. (b) S\&P~500 out-of-sample from 2000: a tight bundle,
entirely above zero. The vertical scales differ; the comparison of interest is
the width of each fan relative to its own level, quantified in
Table~\ref{tab:main}.}
\label{fig:cmp}
\end{figure}

\subsection{What separates them}

A median out-of-sample Sharpe of $0.33$ for plain trend following on an equity
index is unremarkable, and should be: it is approximately the modest,
well-documented figure the literature would lead one to expect. The level is not
the finding.

The finding is the shape. Dispersion falls from $0.16$ to $0.09$. The fraction
of positive constructions rises from $5\%$ to $100\%$. And $\Delta$ falls from
$0.45$ to $0.10$: on the synthetic null, choosing the best construction after
the fact buys $0.45$ of Sharpe that does not exist, whereas on the S\&P it buys
$0.10$, because there is far less to choose \emph{between}.

Figure~\ref{fig:disp} shows the two distributions directly. This is the
specification curve of \cite{simonsohn2020} in histogram form, and it is what we
propose reporting in place of a single equity curve.

\begin{figure}[t]
\centering
\includegraphics[width=\textwidth]{fig/fig3_dispersion.pdf}
\caption{Distribution of annualized out-of-sample Sharpe across the 432
constructions, one observation per construction. Dashed lines mark the medians.
The separation between the two distributions --- in location, in width, and in
support --- is the diagnostic.}
\label{fig:disp}
\end{figure}

\subsection{A caveat on reading the fans directly}

Figure~\ref{fig:cmp} is drawn in cumulative-return space, where the visual
width of a fan is not a measure of construction sensitivity. The S\&P bundle is
in fact \emph{wider} in raw terms --- the 10--90\% band spans $65.6$ percentage
points at the terminal date against $50.2$ for the synthetic null --- simply
because it is compounding a positive return. Sensitivity is a ratio, not a
distance.

Figure~\ref{fig:band} therefore plots each path as a deviation from its
bundle's median path, expressed in multiples of that bundle's median terminal
level. On that scale the terminal 10--90\% band is $1.55\times$ the level for
the synthetic null and $0.78\times$ for the S\&P: twice as tight. This is the
representation we recommend for visual reporting, and it agrees with the
Sharpe dispersion in Table~\ref{tab:main}, as it must.

\begin{figure}[t]
\centering
\includegraphics[width=\textwidth]{fig/fig4_dispersion_band.pdf}
\caption{Construction sensitivity on a common scale. Each faint line is one
construction's deviation from the median construction, divided by the bundle's
median terminal level. Shaded bands are the inter-decile and inter-quartile
ranges. A wide fan (a) means the reported result depends on which construction
was chosen; a tight one (b) means it does not.}
\label{fig:band}
\end{figure}

\subsection{How large must the grid be?}

If reporting a distribution over constructions is the remedy, the choice of
which constructions to enumerate is itself a researcher degree of freedom, and
the objection of infinite regress is fair. It does not, however, bite in
practice, for three reasons.

First, the axes are the protocol's own parameters, not the model's. They are
few, they are enumerable, and each admits a handful of defensible values ---
unlike a model's parameter space, which is unbounded.

Second, the response surface over those axes is smooth. Inserting a fifth
in-sample window length between 24 and 36 months tells one almost nothing that
24 and 36 did not.

Third, and decisively, the required grid size is an empirical question with a
measurable answer. Table~\ref{tab:conv} reports the accuracy with which a
random subgrid of size $n$ recovers the median of the full 432-point grid.

\begin{table}[h]
\centering\small
\caption{Mean absolute error in the estimated grid median, over 4{,}000 random
subgrids of each size, against the full 432-point grid.}
\label{tab:conv}
\begin{tabular}{rrr}
\toprule
Subgrid size $n$ & Synthetic & S\&P 500 \\
\midrule
  5 & 0.066 & 0.054 \\
 10 & 0.045 & 0.040 \\
 20 & 0.032 & 0.034 \\
 50 & 0.020 & 0.025 \\
100 & 0.013 & 0.019 \\
200 & 0.007 & 0.013 \\
\bottomrule
\end{tabular}
\end{table}

The grid needs to be large enough that the estimation error is small relative
to the effect being examined, and no larger. The difference under test here ---
a median of $+0.33$ against $-0.27$, a gap of $0.60$ --- is roughly eighteen
times the estimation error at $n = 20$. Enumerating all 432 buys precision that
the decision does not consume. We report the full grid because it is cheap
(Section~5), not because the conclusion requires it.

The regress terminates, then, not because the grid choice is objective, but
because it is \emph{declared}. A single curve conceals its search; a bundle
reported with its grid exposes the remaining discretion to the reader, who can
object to a specific omitted axis. That is a weaker guarantee than objectivity
and a much stronger one than convention.

\subsection{Marginal effects}

Because the grid is exhaustive, the marginal effect of each construction choice
is available by conditioning. On the S\&P data three of the four principal axes
are flat: median Sharpe varies by less than $0.01$ across in-sample window
lengths ($0.324$--$0.332$) and by $0.01$ across ranking objectives.

One axis is not flat. Rolling windows give a median of $0.23$; expanding windows
give $0.38$. A researcher who adopted an expanding window as a default, reported
the curve and never swept the alternative would have taken an unmeasured $0.15$
Sharpe difference on the strength of an unexamined convention. Notably, the same
axis has the opposite and much weaker sign on the synthetic null ($-0.25$ versus
$-0.30$), which is what one expects of an effect that interacts with the
presence of real signal rather than reflecting a universal preference.

This is the practical case for the bundle, independent of the epistemic one: it
converts every unexamined default into a measurement.

\section{Making exhaustive evaluation cheap}

Exhaustive enumeration is a methodological prescription only if the enumeration
is affordable. We report the optimisation work because its lesson generalises
and because the naive result was misleading.

The workload has an unusual structure: it is \emph{redundant} before it is
\emph{parallel}. Every construction traverses the same candidate matrix, and
constructions differing only in $\rho$ or $K$ produce identical in-sample
rankings, since those parameters act after the ranking.

Instrumented timings on a grid of 576 constructions over a $5000 \times 1000$
candidate matrix attributed $82.9\%$ of runtime to a single operation: the
in-sample ranking objective, recomputed over the full window slice once per
window per construction. Out-of-sample leg assembly accounted for $10.0\%$,
greedy de-correlation for $2.9\%$, and slicing for $1.2\%$. By Amdahl's law
\cite{amdahl1967}, no amount of work on the remaining $17.1\%$ could have
produced more than a $1.2\times$ speedup --- a ceiling we established only after
implementing two optimisations that fell below it.

The effective change was a memoization of the ranking objective keyed on
$(\text{objective}, \text{window})$, shared across constructions within a
process. On the benchmark grid this reduced $5{,}952$ evaluations to $656$, and
the predicted and observed runtimes agree closely:
\[
45.3\,\text{s} \times \left(0.171 + \tfrac{0.829}{9.1}\right) = 11.9\,\text{s}
\quad\text{predicted},\qquad 12.0\,\text{s}\ \text{measured}.
\]

Two results are worth recording for practitioners. First, a closed-form
prefix-sum evaluation of moment-based objectives --- asymptotically the
strongest optimisation available, reducing per-window cost from $O(Tk)$ to
$O(k)$ --- yielded no measurable gain on a wide grid, because the memoization
had already eliminated the repetition it exploits; it returns a $1.5\times$
gain only in the regime where each schedule appears once. Second, the obvious
parallelisation, dispatching one construction per worker, was $5.9\times$
\emph{slower} than serial execution: each task shipped a fresh copy of the
candidate matrix and began with an empty cache, so the eliminated work returned,
multiplied by the number of tasks. Grouping constructions that share a schedule
into the same task restores the gain. Parallelism that destroys locality is a
regression.

\section{Limitations}

The procedure inherits the causality of its inputs and cannot verify it. If a
candidate's P\&L was generated by a rule fitted on the full history, every
number reported here is contaminated regardless of the walk-forward discipline
applied on top.

A tight bundle establishes insensitivity to construction choice. It does not
establish that the edge will persist, that the candidate universe was not itself
the product of an unreported search, or that the strategy is implementable at
size. It answers one question well and is silent on the others.

The dispersion of $\widehat{S}$ over $\mathcal{C}$ is not a sampling
distribution and admits no direct frequentist interpretation; constructions
share data and are strongly dependent. We report it as a descriptive summary of
sensitivity, not as an interval estimate. Formal inference over specification
spaces --- as developed in \cite{simonsohn2020} --- is compatible with this
procedure but is not attempted here.

Finally, the grid is a modelling choice like any other, and a researcher
determined to produce a tight bundle can do so by narrowing it. The grid must be
reported in full alongside the result, and we regard a bundle reported without
its grid as no more informative than a single curve.

\section{Availability}

A reference implementation is released under the MIT licence at
\begin{center}\small\url{https://github.com/epiphanyalpha/walkforward-bundle}\end{center}
\vspace{-0.4em}
Both experiments in this paper reproduce from a single command, with no data
download:

\begin{quote}\small\ttfamily
pip install walkforward-bundle\\
python -m validation.demo
\end{quote}

The synthetic series is generated deterministically from a fixed seed; the
S\&P~500 results are reproduced by passing a daily close series via
\texttt{-{}-csv}. No market data is redistributed with the package.

\begin{thebibliography}{9}\small

\bibitem{amdahl1967} Amdahl, G.~M. (1967). Validity of the single processor
approach to achieving large scale computing capabilities.
\emph{AFIPS Conference Proceedings}, 30, 483--485.

\bibitem{bailey2014} Bailey, D.~H., \& López de Prado, M. (2014). The Deflated
Sharpe Ratio: Correcting for Selection Bias, Backtest Overfitting and
Non-Normality. \emph{Journal of Portfolio Management}, 40(5), 94--107.

\bibitem{bailey2015} Bailey, D.~H., Borwein, J.~M., López de Prado, M., \&
Zhu, Q.~J. (2015). The Probability of Backtest Overfitting. Working paper,
revised February 2015.

\bibitem{gelman2013} Gelman, A., \& Loken, E. (2013). The garden of forking
paths: Why multiple comparisons can be a problem, even when there is no
``fishing expedition''. Working paper, Department of Statistics, Columbia
University.

\bibitem{harvey2016} Harvey, C.~R., Liu, Y., \& Zhu, H. (2016). $\ldots$ and the
Cross-Section of Expected Returns. \emph{Review of Financial Studies}, 29(1),
5--68.

\bibitem{simonsohn2020} Simonsohn, U., Simmons, J.~P., \& Nelson, L.~D. (2020).
Specification curve analysis. \emph{Nature Human Behaviour}, 4, 1208--1214.

\bibitem{steegen2016} Steegen, S., Tuerlinckx, F., Gelman, A., \& Vanpaemel, W.
(2016). Increasing Transparency Through a Multiverse Analysis.
\emph{Perspectives on Psychological Science}, 11(5), 702--712.

\bibitem{white2000} White, H. (2000). A Reality Check for Data Snooping.
\emph{Econometrica}, 68(5), 1097--1126.

\end{thebibliography}

\end{document}
